% Pass5586: close the characteristic-two rank frontier of the Pass5580
% projectivity-graph incidence design using the published modular structure of
% the natural PSL_2(q) projective-line permutation module.
\subsection{Characteristic-two closure: the projectivity-frame rank is $(q+1)^2/2$}
\label{sec:pass5586-binary-psl2-span}

The binary-rank pattern left explicitly conjectural at the end of
Pass~5585 can be closed for every odd prime power $q$.  The important point is
that this is a theorem about the projectivity-graph design of
Section~\ref{sec:pass5580-reye-psl2}; it is \emph{not} the separate W$(3,q)$
footprint-rank conjecture firewalled in Passes~5358 and~5376.

Let $V=\mathbf F_2[\mathbf P^1(q)]$ be the natural permutation module for
$G=\PSL_2(q)$ and let
\[
 \mathcal A=\operatorname{span}_{\mathbf F_2}\{\rho(g):g\in G\}
 \subseteq\operatorname{End}_{\mathbf F_2}(V).
\]
Vectorizing matrices identifies the rows of the Pass~5580 incidence matrix
with the permutation matrices of one $\PGL_2(q)/\PSL_2(q)$ coset.  Left
multiplication by one fixed coset representative is invertible, so
\[
 \operatorname{rank}_2 M=\dim_{\mathbf F_2}\mathcal A.
\]

The needed modular module structure is classical.  Zavarnitsine's Lemma~7
(
\emph{Subextensions for a permutation $\PSL_2(q)$-module}, 2013), refined in
Revin--Zavarnitsine (2021), gives the Loewy shape after scalar extension to
$k=\overline{\mathbf F}_2$:
\[
 \boxed{
 \mathbf 1\ \big|\ (U_+\oplus U_-)\ \big|\ \mathbf 1,
 \qquad
 \dim U_+=\dim U_-=a:=\frac{q-1}{2}.
 }
\]
The bottom trivial module is the all-one vector, the unique maximal submodule
is the augmentation kernel, and the quotient is again trivial.  For
$q\equiv\pm1\pmod8$ the two middle constituents are already absolutely
irreducible over $\mathbf F_2$; for $q\equiv\pm3\pmod8$ the $(q-1)$-dimensional
middle constituent splits into $U_+\oplus U_-$ after extending scalars to
$\mathbf F_4$ (and hence to $k$).

Let $J=\operatorname{rad}(\mathcal A\otimes k)$.  Faithfulness of $V$ for the
image algebra and the three simple composition types give
\[
 \boxed{
 (\mathcal A\otimes k)/J\cong k\oplus M_a(k)\oplus M_a(k),
 }
\]
so
\[
 \dim((\mathcal A\otimes k)/J)=1+2a^2.
\]
The Loewy filtration gives an immediate upper bound on the radical.  Every
element of $J$ kills the one-dimensional socle, maps the $2a$-dimensional
middle layer into the socle, and maps the one-dimensional top into the
middle layer plus socle.  Hence
\[
 \dim J\le 2a+(2a+1)=4a+1.
\]
For the reverse inequality, the two middle simples occur in
$\operatorname{rad}V/\operatorname{rad}^2V$.  Therefore the two distinct Pierce
components of $J/J^2$ from the trivial top to $U_+$ and $U_-$ are nonzero and
contribute at least $a+a$ dimensions.  Transposition preserves the permutation
matrix algebra and its radical and supplies the two reverse Pierce components,
contributing another $a+a$.  Finally
$\operatorname{rad}^2V=\operatorname{soc}V\ne0$, so $J^2\ne0$.  Thus
\[
 \dim J\ge4a+1,
\]
and equality holds.

Consequently
\[
\begin{aligned}
 \operatorname{rank}_2M
 &=\dim\mathcal A\\
 &=1+2a^2+4a+1\\
 &=2(a+1)^2\\
 &=\boxed{\frac{(q+1)^2}{2}}.
\end{aligned}
\]
This proves the all-odd prime-power law.  The previously measured anchors
\[
 q=3,5,7,11,13
 \quad\longmapsto\quad
 8,18,32,72,98
\]
are exact instances rather than evidence for an open formula.

\paragraph{Evidence boundary.}
The theorem closes only the characteristic-two rank of the
$\PSL_2(q)$ projectivity-graph incidence family.  It does not prove the
independent all-odd W$(3,q)$ footprint-rank law, and it supplies no $q>3$
polytope or physical identification.  The older Pass~5358 warning remains
essential: the proof uses the nonsplit modular Loewy structure itself, not an
invalid reduction of the characteristic-zero decomposition.
